\documentclass[12pt,a4paper]{article}
\usepackage[UTF8]{ctex}
\usepackage{geometry}
\usepackage{listings}
\usepackage{xcolor}
\usepackage{hyperref}
\usepackage{booktabs}
\usepackage{float}
\usepackage{tcolorbox}
\usepackage{enumitem}
\usepackage{graphicx}

\geometry{left=2.5cm,right=2.5cm,top=2.5cm,bottom=2.5cm}

\lstset{
    language=Python,
    basicstyle=\ttfamily\small,
    keywordstyle=\color{blue}\bfseries,
    commentstyle=\color{gray},
    stringstyle=\color{orange},
    numbers=left,
    numberstyle=\tiny\color{gray},
    stepnumber=1,
    numbersep=8pt,
    backgroundcolor=\color{gray!10},
    frame=single,
    frameround=fttt,
    breaklines=true,
    showstringspaces=false,
    tabsize=4,
    captionpos=b,
    xleftmargin=2em,
    xrightmargin=1em,
    aboveskip=1em,
    belowskip=1em
}

\tcbuselibrary{skins,breakable}
\newtcolorbox{knowledgebox}[1][]{
    colback=blue!5!white,
    colframe=blue!50!black,
    fonttitle=\bfseries,
    title=#1,
    breakable
}

\newtcolorbox{practicebox}[1][]{
    colback=green!5!white,
    colframe=green!50!black,
    fonttitle=\bfseries,
    title=#1,
    breakable
}

\newtcolorbox{tipbox}[1][]{
    colback=yellow!10!white,
    colframe=orange!70!black,
    fonttitle=\bfseries,
    title=#1,
    breakable
}

\newtcolorbox{theorybox}[1][]{
    colback=purple!5!white,
    colframe=purple!70!black,
    fonttitle=\bfseries,
    title=#1,
    breakable
}

\hypersetup{
    colorlinks=true,
    linkcolor=blue,
    filecolor=magenta,      
    urlcolor=cyan,
    bookmarks=true,
    pdfstartview=FitH
}

\title{\textbf{Pandas数据处理实战手册}\\[0.5em]\large 从零基础到竞赛实战}
\author{竞赛培训组}
\date{\today}

\begin{document}

\maketitle

\begin{tipbox}[学习指引]
\textbf{后续课程}：学完 Pandas 后，下一步是部署大数据平台。

\textbf{配套练习}：练习数据已包含在本手册的 CSV 文件中。
\end{tipbox}

\tableofcontents
\newpage

%========================================
\section{前言：为什么要学习Pandas？}
%========================================

\begin{theorybox}[什么是Pandas？]
\textbf{Pandas}是Python语言中最流行的数据处理和分析库，它的名字来源于"Panel Data"（面板数据）的缩写。Pandas提供了高效、灵活的数据结构，特别适合处理结构化数据（如表格数据）。

\textbf{核心特点：}
\begin{itemize}
    \item \textbf{DataFrame}：二维表格数据结构，类似Excel表格
    \item \textbf{Series}：一维数组数据结构，类似表格中的一列
    \item 强大的数据读取、清洗、转换、分析功能
    \item 与NumPy、Matplotlib等库无缝集成
\end{itemize}
\end{theorybox}

\begin{theorybox}[为什么数据处理很重要？]
在真实的数据分析项目中，数据科学家通常会将\textbf{70\%-80\%}的时间花在数据清洗和预处理上，只有20\%-30\%的时间用于建模和分析。

\textbf{常见数据质量问题：}
\begin{enumerate}
    \item \textbf{缺失值}：某些记录缺少某些字段的值
    \item \textbf{重复值}：同一条数据被记录了多次
    \item \textbf{异常值}：明显不合理的数据（如年龄为负数）
    \item \textbf{格式不一致}：同一类数据有多种表示方式
    \item \textbf{数据类型错误}：数字被存储为字符串
\end{enumerate}

\textbf{竞赛评分标准：}在大数据竞赛中，数据预处理通常占\textbf{30\%}的评分权重，是获得高分的关键！
\end{theorybox}

%========================================
\section{四个练习数据文件介绍}
%========================================

\subsection{数据文件概览}

本手册所有示例代码都围绕以下\textbf{四个CSV文件}展开，请确保这些文件和代码在同一目录下：

\begin{table}[H]
\centering
\caption{练习数据文件说明}
\begin{tabular}{|l|l|l|l|}
\hline
\textbf{文件名} & \textbf{数据量} & \textbf{主题} & \textbf{主要问题} \\
\hline
01\_电商订单数据\_练习.csv & 50行 & 电商订单 & 重复、缺失、异常值、格式混乱 \\
\hline
02\_传感器数据\_练习.csv & 144行 & 传感器读数 & 时间序列、格式问题 \\
\hline
03\_缺失值处理\_练习.csv & 30行 & 员工信息 & 各种缺失值场景 \\
\hline
04\_数据格式清洗\_练习.csv & 30行 & 用户信息 & 格式不统一、单位混乱 \\
\hline
\end{tabular}
\end{table}

\subsection{01\_电商订单数据\_练习.csv}

\begin{knowledgebox}[数据字段说明]
\begin{itemize}
    \item \textbf{order\_id}：订单编号
    \item \textbf{user\_id}：用户ID（部分缺失）
    \item \textbf{product\_name}：商品名称
    \item \textbf{price}：价格（带"元"单位）
    \item \textbf{quantity}：数量（有负数异常值）
    \item \textbf{order\_date}：订单日期（多种格式）
    \item \textbf{address}：地址
    \item \textbf{phone}：手机号（带横线）
    \item \textbf{tags}：商品标签（JSON格式）
\end{itemize}
\end{knowledgebox}

\subsection{02\_传感器数据\_练习.csv}

\begin{knowledgebox}[数据字段说明]
\begin{itemize}
    \item \textbf{timestamp}：时间戳
    \item \textbf{sensor\_id}：传感器ID
    \item \textbf{temperature}：温度
    \item \textbf{humidity}：湿度（部分缺失）
    \item \textbf{pressure}：压力
    \item \textbf{status}：设备状态
\end{itemize}
\end{knowledgebox}

\subsection{03\_缺失值处理\_练习.csv}

\begin{knowledgebox}[数据字段说明]
\begin{itemize}
    \item \textbf{id}：编号
    \item \textbf{name}：姓名
    \item \textbf{age}：年龄（部分缺失）
    \item \textbf{gender}：性别（部分缺失）
    \item \textbf{salary}：工资（部分缺失）
    \item \textbf{department}：部门（部分缺失）
    \item \textbf{join\_date}：入职日期（部分缺失）
    \item \textbf{phone}：电话
    \item \textbf{email}：邮箱
    \item \textbf{score}：评分（部分缺失）
\end{itemize}
\end{knowledgebox}

\subsection{04\_数据格式清洗\_练习.csv}

\begin{knowledgebox}[数据字段说明]
\begin{itemize}
    \item \textbf{id}：编号
    \item \textbf{name}：姓名（带空格）
    \item \textbf{birth\_date}：出生日期（多种格式）
    \item \textbf{height}：身高（带单位）
    \item \textbf{weight}：体重（带单位）
    \item \textbf{income}：收入（带k单位）
    \item \textbf{phone}：电话
    \item \textbf{address}：地址
    \item \textbf{product\_tags}：商品标签
\end{itemize}
\end{knowledgebox}

%========================================
\section{Python与Pandas基础}
%========================================

\subsection{安装与导入}

\begin{theorybox}[什么是Python包？]
Python的强大之处在于它有丰富的第三方库（包）。这些库是由全球开发者贡献的，可以帮助我们快速实现各种功能，而不需要从零开始编写代码。

\textbf{常用数据科学包：}
\begin{itemize}
    \item \textbf{pandas}：数据处理和分析
    \item \textbf{numpy}：数值计算
    \item \textbf{matplotlib}：数据可视化
    \item \textbf{scikit-learn}：机器学习
\end{itemize}
\end{theorybox}

\begin{lstlisting}[caption={安装Pandas}]
# 在命令行中运行
pip install pandas requests
\end{lstlisting}

\begin{lstlisting}[caption={导入Pandas}]
import pandas as pd
import numpy as np
\end{lstlisting}

\begin{tipbox}[导入别名约定]
在Python数据科学社区中，有一些约定俗成的导入别名：
\begin{itemize}
    \item \texttt{import pandas as pd} —— pd代表pandas
    \item \texttt{import numpy as np} —— np代表numpy
\end{itemize}
使用别名可以让代码更简洁，这也是行业内的标准做法。
\end{tipbox}

\subsection{读取CSV文件}

\begin{theorybox}[什么是CSV文件？]
\textbf{CSV}（Comma-Separated Values，逗号分隔值）是一种通用的、简单的数据存储格式。

\textbf{CSV的特点：}
\begin{itemize}
    \item 纯文本格式，可以用任何文本编辑器打开
    \item 每行是一条记录，字段之间用逗号分隔
    \item 第一行通常是列名（表头）
    \item 可以被Excel直接打开和编辑
    \item 文件体积小，兼容性好
\end{itemize}

\textbf{示例CSV内容：}
\begin{verbatim}
name,age,city
张三,25,北京
李四,30,上海
王五,28,广州
\end{verbatim}
\end{theorybox}

\begin{knowledgebox}[知识点：pd.read\_csv()]
\texttt{pd.read\_csv()}是Pandas中最常用的函数，用于读取CSV文件并转换为DataFrame对象。

\textbf{基本语法：}
\begin{verbatim}
df = pd.read_csv('文件路径.csv')
\end{verbatim}

\textbf{常用参数：}
\begin{itemize}
    \item \texttt{encoding}：指定文件编码，如'utf-8'、'gbk'
    \item \texttt{header}：指定哪一行作为列名，默认为0（第一行）
    \item \texttt{sep}：指定分隔符，默认为逗号
\end{itemize}
\end{knowledgebox}

\begin{lstlisting}[caption={读取四个练习文件}]
# 读取电商订单数据
df_order = pd.read_csv('01_电商订单数据_练习.csv')

# 读取传感器数据
df_sensor = pd.read_csv('02_传感器数据_练习.csv')

# 读取缺失值练习数据
df_missing = pd.read_csv('03_缺失值处理_练习.csv')

# 读取格式清洗练习数据
df_messy = pd.read_csv('04_数据格式清洗_练习.csv')
\end{lstlisting}

\begin{practicebox}[动手练习1：读取并查看数据]
\textbf{任务：}读取四个CSV文件，查看每个文件的基本信息。

\begin{lstlisting}
import pandas as pd

# 读取四个文件
df_order = pd.read_csv('01_电商订单数据_练习.csv')
df_sensor = pd.read_csv('02_传感器数据_练习.csv')
df_missing = pd.read_csv('03_缺失值处理_练习.csv')
df_messy = pd.read_csv('04_数据格式清洗_练习.csv')

# 查看电商订单数据的形状
print("电商订单数据形状：", df_order.shape)
print("\n前5行：")
print(df_order.head())
\end{lstlisting}
\end{practicebox}

%========================================
\section{数据查看与探索}
%========================================

\subsection{查看数据基本信息}

\begin{theorybox}[为什么需要先查看数据？]
在进行任何数据处理之前，我们必须先\textbf{了解数据}。就像医生看病需要先诊断一样，数据分析也需要先"诊断"数据的质量和结构。

\textbf{数据探索的目的：}
\begin{enumerate}
    \item 了解数据规模（多少行、多少列）
    \item 了解每列的数据类型
    \item 发现明显的数据质量问题
    \item 为后续的数据清洗制定策略
\end{enumerate}
\end{theorybox}

\begin{lstlisting}[caption={查看数据形状和列名}]
# 查看数据形状（行数，列数）
print(df_order.shape)

# 查看列名
print(df_order.columns)

# 查看数据类型
print(df_order.dtypes)
\end{lstlisting}

\begin{knowledgebox}[shape属性的含义]
\texttt{df.shape}返回一个元组\texttt{(行数, 列数)}，表示数据的维度。

\textbf{示例：}
\begin{verbatim}
>>> df.shape
(50, 9)  # 表示50行，9列
\end{verbatim}
\end{knowledgebox}

\subsection{查看数据内容}

\begin{lstlisting}[caption={查看数据内容}]
# 查看前5行
print(df_order.head())

# 查看后5行
print(df_order.tail())

# 查看前10行
print(df_order.head(10))
\end{lstlisting}

\begin{tipbox}[head()和tail()的用法]
\begin{itemize}
    \item \texttt{head(n)}：查看前n行，默认n=5
    \item \texttt{tail(n)}：查看后n行，默认n=5
\end{itemize}
查看数据的开头可以了解数据结构，查看结尾可以检查数据是否完整导入。
\end{tipbox}

\subsection{统计信息}

\begin{theorybox}[描述性统计]
描述性统计是对数据特征进行总结和描述的方法。通过统计量，我们可以快速了解数据的分布情况。

\textbf{常用统计量：}
\begin{itemize}
    \item \textbf{count}：非空值数量
    \item \textbf{mean}：平均值
    \item \textbf{std}：标准差（数据离散程度）
    \item \textbf{min}：最小值
    \item \textbf{25\%}：第一四分位数（Q1）
    \item \textbf{50\%}：中位数
    \item \textbf{75\%}：第三四分位数（Q3）
    \item \textbf{max}：最大值
\end{itemize}
\end{theorybox}

\begin{lstlisting}[caption={使用describe()查看统计信息}]
# 查看数值列的统计信息
print(df_order.describe())

# 查看缺失值统计
print(df_order.isnull().sum())
\end{lstlisting}

\begin{knowledgebox}[知识点：isnull()和sum()]
\texttt{isnull()}函数用于检测缺失值，返回一个布尔值的DataFrame，缺失值位置为True。

\texttt{sum()}函数对布尔值进行求和时，True被当作1，False被当作0，因此可以直接统计每列的缺失值数量。

\textbf{执行流程：}
\begin{enumerate}
    \item \texttt{df.isnull()} → 生成布尔值表格
    \item \texttt{.sum()} → 对每列求和（统计True的数量）
\end{enumerate}
\end{knowledgebox}

\begin{practicebox}[动手练习2：探索传感器数据]
\textbf{任务：}使用02\_传感器数据\_练习.csv，查看数据的基本统计信息。

\begin{lstlisting}
import pandas as pd

df_sensor = pd.read_csv('02_传感器数据_练习.csv')

# 查看形状
print("数据形状：", df_sensor.shape)

# 查看统计信息
print("\n统计摘要：")
print(df_sensor.describe())

# 查看有多少个不同的传感器
print("\n传感器数量：", df_sensor['sensor_id'].nunique())
print("传感器列表：", df_sensor['sensor_id'].unique())
\end{lstlisting}
\end{practicebox}

%========================================
\begin{practicebox}[平台实战：用仿真平台数据练习查看与探索]
\begin{lstlisting}[language=Python]
import requests, pandas as pd

df = pd.DataFrame(requests.get("http://163.7.1.176:8889/api/data?n=1000").json())

# 练习：查看数据基本信息
print(df.shape)         # 多少行多少列？
print(df.dtypes)        # 每列什么类型？
print(df.isnull().sum())  # 每列多少缺失值？

# 练习：查看数据内容
print(df.head(10))      # 前10行
print(df.tail(5))       # 后5行

# 练习：统计信息
print(df['price'].describe())  # 价格的统计摘要
print(df['event_type'].value_counts())  # 各事件类型数量
\end{lstlisting}
\end{practicebox}

\section{数据选择}
%========================================

\subsection{选择列}

\begin{theorybox}[DataFrame的列选择]
在Pandas中，选择列是最常用的操作之一。DataFrame可以看作是由多个Series（列）组成的字典。

\textbf{两种选择方式：}
\begin{enumerate}
    \item \textbf{单列选择}：返回一个Series对象
    \item \textbf{多列选择}：返回一个DataFrame对象
\end{enumerate}
\end{theorybox}

\begin{lstlisting}[caption={选择单列和多列}]
# 选择单列（商品名称）
products = df_order['product_name']

# 选择多列（订单ID、商品名称、价格）
order_info = df_order[['order_id', 'product_name', 'price']]
\end{lstlisting}

\begin{tipbox}[注意方括号的区别]
\begin{itemize}
    \item \texttt{df['列名']} —— 单层方括号，选择单列，返回Series
    \item \texttt{df[['列1', '列2']]} —— 双层方括号，选择多列，返回DataFrame
\end{itemize}
\end{tipbox}

\subsection{选择行}

\begin{theorybox}[iloc和loc的区别]
Pandas提供了两种主要的行选择方式：

\begin{enumerate}
    \item \textbf{iloc}：基于\textbf{整数位置}索引（从0开始）
    \item \textbf{loc}：基于\textbf{标签}索引（使用行标签）
\end{enumerate}

对于刚读取的CSV文件，行标签默认就是整数0, 1, 2...，所以iloc和loc看起来效果相似。但当数据经过筛选、排序后，行标签可能不再是连续的整数，这时两者的区别就会显现。
\end{theorybox}

\begin{lstlisting}[caption={使用iloc选择行}]
# 选择第1行（索引为0）
row = df_order.iloc[0]

# 选择前5行
rows = df_order.iloc[0:5]

# 选择第10到20行
rows = df_order.iloc[10:20]
\end{lstlisting}

\subsection{条件筛选}

\begin{theorybox}[布尔索引]
条件筛选是数据分析中最强大的功能之一。它的原理是\textbf{布尔索引}：

\begin{enumerate}
    \item 首先创建一个布尔条件（如\texttt{df['age'] > 18}）
    \item 这个条件会生成一个布尔值的Series
    \item 将这个布尔Series放入方括号中，Pandas会只返回True对应的行
\end{enumerate}

\textbf{常用比较运算符：}
\begin{itemize}
    \item \texttt{==} 等于
    \item \texttt{!=} 不等于
    \item \texttt{>} 大于
    \item \texttt{<} 小于
    \item \texttt{>=} 大于等于
    \item \texttt{<=} 小于等于
\end{itemize}
\end{theorybox}

\begin{lstlisting}[caption={条件筛选示例}]
# 筛选数量大于1的订单
large_orders = df_order[df_order['quantity'] > 1]

# 筛选特定商品
iphone_orders = df_order[df_order['product_name'] == 'iPhone 14 Pro']
\end{lstlisting}

\begin{practicebox}[动手练习3：筛选缺失值数据]
\textbf{任务：}使用03\_缺失值处理\_练习.csv，筛选出年龄缺失的员工。

\begin{lstlisting}
import pandas as pd

df_missing = pd.read_csv('03_缺失值处理_练习.csv')

# 筛选年龄缺失的行
missing_age = df_missing[df_missing['age'].isnull()]

print("年龄缺失的员工：")
print(missing_age[['name', 'age', 'department']])

print(f"\n共 {len(missing_age)} 人年龄缺失")
\end{lstlisting}
\end{practicebox}

%========================================
\begin{practicebox}[平台实战：用仿真平台数据练习选择]
\begin{lstlisting}[language=Python]
# 继续使用上面获取的 df

# 练习：选择列
print(df[['user_id', 'event_type', 'price']].head())

# 练习：条件筛选
buy_data = df[df['event_type'] == 'buy']
print(f"购买记录: {len(buy_data)} 条")

# 练习：组合条件
expensive = df[(df['price'] > 5000) & (df['event_type'] == 'view')]
print(f"高价浏览: {len(expensive)} 条")
\end{lstlisting}
\end{practicebox}

\section{数据清洗：处理缺失值}
%========================================

\subsection{什么是缺失值？}

\begin{theorybox}[缺失值的产生原因]
缺失值（Missing Values）是指数据集中某些记录在某些字段上没有值。缺失值在真实数据中非常常见。

\textbf{缺失值的常见原因：}
\begin{enumerate}
    \item \textbf{数据采集失败}：设备故障、网络中断
    \item \textbf{数据录入遗漏}：人工录入时漏填
    \item \textbf{信息不适用}：某些字段对某些记录不适用
    \item \textbf{隐私保护}：敏感信息被故意隐藏
    \item \textbf{数据合并问题}：多个数据源合并时字段不匹配
\end{enumerate}

\textbf{缺失值的表示方式：}
\begin{itemize}
    \item Python中：\texttt{None}或\texttt{np.nan}（Not a Number）
    \item CSV文件中：空单元格或特定标记（如"N/A"、"NULL"）
    \item 数据库中：\texttt{NULL}
\end{itemize}
\end{theorybox}

\subsection{检测缺失值}

\begin{lstlisting}[caption={检测缺失值}]
# 统计每列的缺失值数量
print(df_missing.isnull().sum())

# 计算缺失值比例
missing_ratio = df_missing.isnull().sum() / len(df_missing) * 100
print(missing_ratio)
\end{lstlisting}

\begin{knowledgebox}[缺失值检测方法]
\begin{itemize}
    \item \texttt{df.isnull()}：返回布尔值DataFrame，缺失值为True
    \item \texttt{df.notnull()}：返回布尔值DataFrame，非缺失值为True
    \item \texttt{df.isnull().sum()}：统计每列缺失值数量
    \item \texttt{df.isnull().sum().sum()}：统计整个DataFrame的缺失值总数
\end{itemize}
\end{knowledgebox}

\subsection{缺失值处理策略}

\begin{theorybox}[缺失值处理策略]
处理缺失值没有绝对正确的方法，需要根据具体情况选择策略：

\textbf{策略一：删除缺失值}
\begin{itemize}
    \item 适用场景：缺失值比例很小（如$<$5\%），且缺失是随机的
    \item 优点：简单直接
    \item 缺点：可能丢失有价值的信息
\end{itemize}

\textbf{策略二：填充缺失值}
\begin{itemize}
    \item 适用场景：缺失值比例较大，或数据很珍贵不能删除
    \item 填充方法：
    \begin{itemize}
        \item 固定值填充（如0、"未知"）
        \item 统计值填充（均值、中位数、众数）
        \item 前后值填充（时间序列数据）
        \item 模型预测填充（高级方法）
    \end{itemize}
\end{itemize}

\textbf{策略三：标记缺失}
\begin{itemize}
    \item 创建一个新列标记哪些值是填充的
    \item 适用于缺失本身可能有意义的情况
\end{itemize}
\end{theorybox}

\subsection{填充缺失值}

\begin{lstlisting}[caption={填充缺失值的方法}]
# 用固定值填充
df_missing['gender'] = df_missing['gender'].fillna('未知')

# 用平均值填充（工资列）
mean_salary = df_missing['salary'].mean()
df_missing['salary'] = df_missing['salary'].fillna(mean_salary)

# 用中位数填充（年龄列）
median_age = df_missing['age'].median()
df_missing['age'] = df_missing['age'].fillna(median_age)

# 用众数填充（部门列）
mode_dept = df_missing['department'].mode()[0]
df_missing['department'] = df_missing['department'].fillna(mode_dept)
\end{lstlisting}

\begin{theorybox}[均值、中位数、众数的选择]
\begin{itemize}
    \item \textbf{均值（Mean）}：适用于数据分布对称的情况，受异常值影响大
    \item \textbf{中位数（Median）}：适用于有异常值或分布偏斜的情况，更稳健
    \item \textbf{众数（Mode）}：适用于分类数据，表示出现频率最高的值
\end{itemize}
\end{theorybox}

\begin{practicebox}[动手练习4：处理缺失值]
\textbf{任务：}使用03\_缺失值处理\_练习.csv，完成缺失值处理。

\begin{lstlisting}
import pandas as pd

df = pd.read_csv('03_缺失值处理_练习.csv')

print("处理前缺失值统计：")
print(df.isnull().sum())

# 用中位数填充年龄
median_age = df['age'].median()
df['age'] = df['age'].fillna(median_age)

# 用众数填充性别
mode_gender = df['gender'].mode()[0]
df['gender'] = df['gender'].fillna(mode_gender)

# 用平均值填充工资
mean_salary = df['salary'].mean()
df['salary'] = df['salary'].fillna(mean_salary)

print("\n处理后缺失值统计：")
print(df.isnull().sum())
\end{lstlisting}
\end{practicebox}

%========================================
\section{数据清洗：处理重复值}
%========================================

\subsection{什么是重复值？}

\begin{theorybox}[重复值的产生原因]
重复值是指数据集中完全相同的记录出现了多次。

\textbf{重复值的常见原因：}
\begin{enumerate}
    \item \textbf{数据录入错误}：同一条数据被重复录入
    \item \textbf{数据合并}：多个数据源合并时产生重复
    \item \textbf{系统问题}：网络重传、事务重复提交
    \item \textbf{定时任务}：数据同步任务重复执行
\end{enumerate}

\textbf{重复值的影响：}
\begin{itemize}
    \item 导致统计分析结果偏差
    \item 增加存储和计算成本
    \item 影响机器学习模型训练
\end{itemize}
\end{theorybox}

\subsection{检测和删除重复值}

\begin{lstlisting}[caption={处理重复值}]
# 检测重复值
duplicates = df_order.duplicated()
print(f"重复行数：{duplicates.sum()}")

# 查看重复的行
duplicate_rows = df_order[df_order.duplicated(keep=False)]
print(duplicate_rows)

# 删除重复值（保留第一个）
df_order_clean = df_order.drop_duplicates()
\end{lstlisting}

\begin{knowledgebox}[duplicated()和drop\_duplicates()详解]
\texttt{duplicated()}函数：
\begin{itemize}
    \item 返回布尔值Series，标记每行是否重复
    \item \texttt{keep='first'}：第一次出现标记为False，后续重复标记为True
    \item \texttt{keep='last'}：最后一次出现标记为False，前面重复标记为True
    \item \texttt{keep=False}：所有重复行都标记为True
\end{itemize}

\texttt{drop\_duplicates()}函数：
\begin{itemize}
    \item 删除重复行，返回新的DataFrame
    \item \texttt{keep='first'}：保留第一次出现的行（默认）
    \item \texttt{keep='last'}：保留最后一次出现的行
    \item \texttt{keep=False}：删除所有重复行
\end{itemize}
\end{knowledgebox}

\begin{practicebox}[动手练习5：处理电商订单重复值]
\textbf{任务：}使用01\_电商订单数据\_练习.csv，检测并删除重复订单。

\begin{lstlisting}
import pandas as pd

df = pd.read_csv('01_电商订单数据_练习.csv')

print(f"原始数据行数：{len(df)}")

# 检测重复值
n_duplicates = df.duplicated().sum()
print(f"重复行数：{n_duplicates}")

# 查看重复的行
duplicate_rows = df[df.duplicated(keep=False)]
print("\n重复的行：")
print(duplicate_rows[['order_id', 'product_name', 'price']])

# 删除重复值
df_clean = df.drop_duplicates()
print(f"\n删除后行数：{len(df_clean)}")
\end{lstlisting}
\end{practicebox}

%========================================
\section{数据清洗：处理异常值}
%========================================

\subsection{什么是异常值？}

\begin{theorybox}[异常值的定义和类型]
异常值（Outliers）是指与数据集中大多数观测值明显不同的数据点。

\textbf{异常值的类型：}
\begin{enumerate}
    \item \textbf{错误值}：数据采集或录入错误（如年龄为负数）
    \item \textbf{真实异常}：真实存在但罕见的情况（如收入极高的人）
    \item \textbf{特殊情况}：特定条件下的正常值（如促销期间的销量激增）
\end{enumerate}

\textbf{异常值的影响：}
\begin{itemize}
    \item 使均值、标准差等统计量失真
    \item 影响机器学习模型的训练效果
    \item 可能导致错误的业务决策
\end{itemize}
\end{theorybox}

\subsection{异常值检测方法}

\begin{theorybox}[IQR方法（箱线图法）]
\textbf{IQR}（Interquartile Range，四分位距）是一种常用的异常值检测方法。

\textbf{计算步骤：}
\begin{enumerate}
    \item 计算第一四分位数 Q1（25\%分位数）
    \item 计算第三四分位数 Q3（75\%分位数）
    \item 计算 IQR = Q3 - Q1
    \item 定义异常值范围：
    \begin{itemize}
        \item 下限 = Q1 - 1.5 × IQR
        \item 上限 = Q3 + 1.5 × IQR
    \end{itemize}
    \item 超出此范围的值被认为是异常值
\end{enumerate}

\textbf{为什么是1.5倍？} 这是统计学中的经验法则，约涵盖99.3\%的正态分布数据。
\end{theorybox}

\begin{lstlisting}[caption={使用IQR方法检测异常值}]
# 计算IQR
Q1 = df_order['quantity'].quantile(0.25)
Q3 = df_order['quantity'].quantile(0.75)
IQR = Q3 - Q1

# 定义异常值范围
lower_bound = Q1 - 1.5 * IQR
upper_bound = Q3 + 1.5 * IQR

# 找出异常值
outliers = df_order[(df_order['quantity'] < lower_bound) | 
                    (df_order['quantity'] > upper_bound)]
print(f"异常值数量：{len(outliers)}")
\end{lstlisting}

\subsection{处理异常值}

\begin{lstlisting}[caption={处理异常值}]
# 将负数数量修正为1
df_order.loc[df_order['quantity'] < 0, 'quantity'] = 1
\end{lstlisting}

\begin{theorybox}[异常值处理策略]
\begin{enumerate}
    \item \textbf{删除}：如果确定是错误值，直接删除
    \item \textbf{修正}：根据业务逻辑修正（如负数数量改为正数）
    \item \textbf{替换}：用边界值或统计值替换
    \item \textbf{保留}：如果是真实异常情况，保留并单独分析
    \item \textbf{分箱}：将异常值归入特殊类别
\end{enumerate}
\end{theorybox}

%========================================
\section{数据格式转换}
%========================================

\subsection{为什么需要格式转换？}

\begin{theorybox}[数据类型的重要性]
正确的数据类型对于数据分析至关重要：

\begin{itemize}
    \item \textbf{数值计算}：价格必须是数值型才能进行加减乘除
    \item \textbf{内存效率}：正确的类型可以节省内存
    \item \textbf{函数兼容性}：某些函数只接受特定类型
    \item \textbf{排序和比较}：字符串"10"小于"2"，但数字10大于2
\end{itemize}

\textbf{常见类型问题：}
\begin{enumerate}
    \item 数字被存储为字符串（如"123元"）
    \item 日期被存储为字符串
    \item 分类数据没有统一格式
\end{enumerate}
\end{theorybox}

\subsection{字符串处理}

\begin{lstlisting}[caption={处理价格列}]
# 查看价格列的格式
print(df_order['price'].head())

# 去除"元"字，转换为数值
df_order['price_num'] = df_order['price'].str.replace('元', '')
df_order['price_num'] = df_order['price_num'].astype(float)
\end{lstlisting}

\begin{knowledgebox}[字符串方法 str.XXX]
Pandas的Series对象有一个\texttt{str}属性，提供了一系列字符串处理方法：

\begin{itemize}
    \item \texttt{str.replace(a, b)}：将a替换为b
    \item \texttt{str.strip()}：去除首尾空格
    \item \texttt{str.lower()} / \texttt{str.upper()}：大小写转换
    \item \texttt{str.contains()}：判断是否包含某子串
    \item \texttt{str.split()}：分割字符串
\end{itemize}
\end{knowledgebox}

\subsection{处理手机号}

\begin{lstlisting}[caption={标准化手机号}]
# 去除横线
df_order['phone_clean'] = df_order['phone'].str.replace('-', '')
\end{lstlisting}

\begin{practicebox}[动手练习6：清洗格式混乱数据]
\textbf{任务：}使用04\_数据格式清洗\_练习.csv，清洗姓名和收入列。

\begin{lstlisting}
import pandas as pd

df = pd.read_csv('04_数据格式清洗_练习.csv')

# 查看原始数据
print("姓名处理前：")
print(df['name'].head())

# 去除姓名中的空格
df['name_clean'] = df['name'].str.strip()

print("\n姓名处理后：")
print(df['name_clean'].head())

# 处理收入（将12k转换为12000）
def convert_income(income_str):
    income_str = str(income_str).strip()
    if 'k' in income_str.lower():
        number = float(income_str.lower().replace('k', ''))
        return number * 1000
    else:
        return float(income_str)

df['income_num'] = df['income'].apply(convert_income)

print("\n收入转换：")
print(df[['income', 'income_num']].head(10))
\end{lstlisting}
\end{practicebox}

%========================================
\section{数据分组与统计}
%========================================

\subsection{分组统计的概念}

\begin{theorybox}[什么是分组统计？]
分组统计（GroupBy）是数据分析中最强大的功能之一，它的思想是：

\textbf{"先分组，再统计"}

\textbf{类比理解：}
\begin{itemize}
    \item 就像老师统计每个班级的平均分
    \item 先按班级\textbf{分组}，再计算每组的\textbf{平均分}
\end{itemize}

\textbf{SQL中的等价概念：}
\begin{verbatim}
SELECT department, AVG(salary) 
FROM employees 
GROUP BY department;
\end{verbatim}
\end{theorybox}

\subsection{分组统计}

\begin{lstlisting}[caption={按部门统计工资}]
# 按部门分组，计算平均工资
dept_avg = df_missing.groupby('department')['salary'].mean()
print(dept_avg)

# 多列分组统计
stats = df_missing.groupby('department').agg({
    'salary': ['mean', 'max', 'min'],
    'age': 'mean'
})
print(stats)
\end{lstlisting}

\begin{knowledgebox}[groupby()详解]
\texttt{groupby()}的基本流程：
\begin{enumerate}
    \item \textbf{Split}：按指定列的值将数据分组
    \item \textbf{Apply}：对每个组应用统计函数
    \item \textbf{Combine}：将结果合并成新的DataFrame
\end{enumerate}

\textbf{常用聚合函数：}
\begin{itemize}
    \item \texttt{mean()}：平均值
    \item \texttt{sum()}：求和
    \item \texttt{count()}：计数
    \item \texttt{max()} / \texttt{min()}：最大/最小值
    \item \texttt{std()}：标准差
\end{itemize}
\end{knowledgebox}

\begin{practicebox}[动手练习7：统计电商数据]
\textbf{任务：}使用01\_电商订单数据\_练习.csv，统计各商品的总销量。

\begin{lstlisting}
import pandas as pd

df = pd.read_csv('01_电商订单数据_练习.csv')

# 处理价格（去除"元"字）
df['price_num'] = df['price'].str.replace('元', '').astype(float)

# 计算订单金额
df['total_amount'] = df['price_num'] * df['quantity']

# 按商品分组统计
product_stats = df.groupby('product_name').agg({
    'quantity': 'sum',
    'total_amount': 'sum'
}).sort_values('total_amount', ascending=False)

print("商品销售统计TOP10：")
print(product_stats.head(10))
\end{lstlisting}
\end{practicebox}

%========================================
\begin{practicebox}[平台实战：用仿真平台数据练习分组统计]
\begin{lstlisting}[language=Python]
# 练习：按品类统计
print(df.groupby('category')['price'].mean().sort_values(ascending=False))

# 练习：按设备统计
print(df['device'].value_counts())

# 练习：按事件类型和品类交叉统计
pivot = df.pivot_table(values='price', index='category',
                       columns='event_type', aggfunc='count', fill_value=0)
print(pivot)
\end{lstlisting}
\end{practicebox}

\section{综合实战：完整数据处理流程}
%========================================

\subsection{数据处理的完整流程}

\begin{theorybox}[标准数据处理流程]
一个完整的数据处理项目通常遵循以下流程：

\begin{enumerate}
    \item \textbf{数据读取}：从文件或数据库加载数据
    \item \textbf{数据探索}：了解数据的基本情况
    \item \textbf{数据清洗}：
    \begin{itemize}
        \item 处理缺失值
        \item 处理重复值
        \item 处理异常值
        \item 格式转换
    \end{itemize}
    \item \textbf{特征工程}：创建新特征、转换现有特征
    \item \textbf{数据分析}：统计分析、可视化
    \item \textbf{结果保存}：保存清洗后的数据
\end{enumerate}
\end{theorybox}

\subsection{案例：完整的电商订单数据处理}

\begin{lstlisting}[caption={完整的数据处理流程}]
import pandas as pd
import numpy as np

# 第1步：读取数据
df = pd.read_csv('01_电商订单数据_练习.csv')
print(f"原始数据：{len(df)}行")

# 第2步：删除重复值
df = df.drop_duplicates()
print(f"删除重复后：{len(df)}行")

# 第3步：处理缺失值
df['user_id'] = df['user_id'].fillna('UNKNOWN')

# 第4步：处理异常值（负数数量）
df.loc[df['quantity'] < 0, 'quantity'] = 1

# 第5步：格式转换（价格）
df['price_num'] = df['price'].str.replace('元', '').astype(float)

# 第6步：格式转换（手机号）
df['phone_clean'] = df['phone'].str.replace('-', '')

# 第7步：计算订单金额
df['total_amount'] = df['price_num'] * df['quantity']

# 第8步：统计分析
product_sales = df.groupby('product_name')['total_amount'].sum()
print("\n各商品销售额：")
print(product_sales.sort_values(ascending=False).head())

# 第9步：保存结果
df.to_csv('清洗后的订单数据.csv', index=False)
print("\n数据已保存到：清洗后的订单数据.csv")
\end{lstlisting}

%========================================
\section{云端仿真平台实战}

\begin{knowledgebox}[这一章要做什么？]
前面的练习用的是本地 CSV 文件。这一章教你从云端数据仿真平台实时获取数据，用 Pandas 做分析。

\textbf{仿真平台地址}：\texttt{http://163.7.1.176:8889}

平台每秒自动生成电商用户行为数据（浏览、点击、加购、购买等），你可以通过 API 获取。
\end{knowledgebox}

\subsection{用 requests 获取数据}

\begin{lstlisting}[language=Python]
import requests
import pandas as pd

# 从仿真平台获取 500 条数据
response = requests.get("http://163.7.1.176:8889/api/data?n=500")
data = response.json()  # 解析 JSON

# 转换为 DataFrame
df = pd.DataFrame(data)

# 查看数据
print(df.shape)         # (500, 11)
print(df.columns)       # 列名
print(df.head())        # 前5行
print(df.dtypes)        # 数据类型
\end{lstlisting}

\subsection{数据清洗练习——脏数据实战}

\begin{practicebox}[练习1：用脏数据练习清洗（模拟比赛场景）]
比赛时数据"可能不完整、不一致、重复、含噪音"。平台提供了脏数据 API：

\begin{lstlisting}[language=Python]
import requests
import pandas as pd

# 获取含脏数据的数据集
dirty = pd.DataFrame(requests.get(
    "http://163.7.1.176:8889/api/dirty?n=500").json())

# 检查脏数据情况
print(dirty.isnull().sum())  # 缺失值统计
print(dirty.duplicated().sum())  # 重复行数
print(dirty['price'].describe())  # 价格统计（看是否有异常值）

# 清洗步骤
dirty_clean = dirty.copy()

# 1. 删除重复行
dirty_clean = dirty_clean.drop_duplicates()
print(f"去重后: {len(dirty_clean)} 行")

# 2. 处理缺失值
dirty_clean['price'] = dirty_clean['price'].fillna(dirty_clean['price'].median())
dirty_clean['category'] = dirty_clean['category'].fillna('未知')
dirty_clean['stay_time'] = dirty_clean['stay_time'].fillna(0)

# 3. 处理异常值（IQR法则）
Q1 = dirty_clean['price'].quantile(0.25)
Q3 = dirty_clean['price'].quantile(0.75)
IQR = Q3 - Q1
dirty_clean['price'] = dirty_clean['price'].clip(
    lower=max(0, Q1 - 1.5*IQR),
    upper=Q3 + 1.5*IQR
)

# 4. 字符串清洗
dirty_clean['category'] = dirty_clean['category'].str.strip()

print(f"清洗后: {len(dirty_clean)} 行")
print(dirty_clean.isnull().sum())  # 确认无缺失值
\end{lstlisting}
\end{practicebox}

\begin{practicebox}[练习2：CSV 格式数据清洗]
\begin{lstlisting}[language=Python]
# 从平台获取 CSV 格式的脏数据
import requests
import pandas as pd

response = requests.get("http://163.7.1.176:8889/api/dirty/csv?n=500")
from io import StringIO
df = pd.read_csv(StringIO(response.text))

# 清洗
df = df.dropna(subset=['category'])  # 删除品类为空的行
df['price'] = pd.to_numeric(df['price'], errors='coerce')  # 强制转数字
df['stay_time'] = pd.to_numeric(df['stay_time'], errors='coerce')

# 导出清洗后的数据
df.to_csv('cleaned.csv', index=False)
print(f"清洗后 {len(df)} 行，已保存为 cleaned.csv")
\end{lstlisting}
\end{practicebox}

\begin{practicebox}[练习3：从仿真平台获取数据并清洗（原始练习）]
\begin{lstlisting}[language=Python]
import requests
import pandas as pd

# 获取 1000 条数据
df = pd.DataFrame(requests.get("http://163.7.1.176:8889/api/data?n=1000").json())

# 练习1：检查缺失值
print(df.isnull().sum())

# 练习2：price 列有缺失值（不是所有事件都有价格）
# 用 0 填充缺失的价格
df['price'] = df['price'].fillna(0)

# 练习3：将 timestamp 转为日期类型
df['timestamp'] = pd.to_datetime(df['timestamp'])

# 练习4：提取小时
df['hour'] = df['timestamp'].dt.hour

# 练习5：按事件类型统计
print(df['event_type'].value_counts())
\end{lstlisting}
\end{practicebox}

\subsection{数据分析练习}

\begin{practicebox}[练习2：分析仿真平台数据]
\begin{lstlisting}[language=Python]
import requests
import pandas as pd

df = pd.DataFrame(requests.get("http://163.7.1.176:8889/api/data?n=2000").json())

# 练习1：各品类的平均价格
print(df.groupby('category')['price'].mean().sort_values(ascending=False))

# 练习2：各设备的使用占比
print(df['device'].value_counts(normalize=True))

# 练习3：购买记录筛选
buys = df[df['event_type'] == 'buy']
print(f"购买记录: {len(buys)} 条")
print(buys.groupby('category')['price'].sum().sort_values(ascending=False))

# 练习4：各来源页面的转化情况
pivot = df.pivot_table(values='user_id', index='page_source',
                       columns='event_type', aggfunc='count', fill_value=0)
print(pivot)
\end{lstlisting}
\end{practicebox}

\subsection{导出清洗后的数据}

\begin{lstlisting}[language=Python]
# 清洗后导出为 CSV（后续大数据分析会用到）
df_clean = df.dropna(subset=['price'])
df_clean.to_csv('cleaned_behavior.csv', index=False, encoding='utf-8-sig')
print(f"已导出 {len(df_clean)} 条清洗后的数据")
\end{lstlisting}

\begin{tipbox}[衔接下一步]
清洗后的数据 \texttt{cleaned\_behavior.csv} 可以上传到大数据平台（HDFS），在《大数据分析实战手册》中用 Spark SQL 进一步分析。
\end{tipbox}

\section{竞赛模拟题}
%========================================

\subsection{题目一：电商订单数据清洗}

\begin{practicebox}[竞赛模拟题1]
\textbf{题目要求：}

使用\textbf{01\_电商订单数据\_练习.csv}，完成以下任务：

\begin{enumerate}
    \item 读取数据，查看基本信息（形状、列名、前5行）
    \item 统计缺失值数量
    \item 删除重复订单
    \item 处理user\_id缺失值（用'UNKNOWN'填充）
    \item 处理负数数量（修正为1）
    \item 提取价格中的数字，转换为数值型
    \item 标准化手机号格式（去除横线）
    \item 计算每个订单的总金额（价格×数量）
    \item 统计每个商品的总销售额和总销量
    \item 保存清洗后的数据到新的CSV文件
\end{enumerate}

\textbf{参考答案：}
\begin{lstlisting}
import pandas as pd

# 读取数据
df = pd.read_csv('01_电商订单数据_练习.csv')

# 查看基本信息
print(f"数据形状：{df.shape}")
print(f"列名：{df.columns.tolist()}")
print(df.head())

# 统计缺失值
print("\n缺失值统计：")
print(df.isnull().sum())

# 删除重复值
df = df.drop_duplicates()

# 填充缺失值
df['user_id'] = df['user_id'].fillna('UNKNOWN')

# 处理异常值
df.loc[df['quantity'] < 0, 'quantity'] = 1

# 格式转换
df['price_num'] = df['price'].str.replace('元', '').astype(float)
df['phone_clean'] = df['phone'].str.replace('-', '')

# 计算金额
df['total_amount'] = df['price_num'] * df['quantity']

# 统计分析
product_stats = df.groupby('product_name').agg({
    'total_amount': 'sum',
    'quantity': 'sum'
}).sort_values('total_amount', ascending=False)

print("\n商品销售统计：")
print(product_stats.head(10))

# 保存结果
df.to_csv('清洗后的订单数据.csv', index=False)
\end{lstlisting}
\end{practicebox}

\subsection{题目二：员工信息数据处理}

\begin{practicebox}[竞赛模拟题2]
\textbf{题目要求：}

使用\textbf{03\_缺失值处理\_练习.csv}，完成以下任务：

\begin{enumerate}
    \item 读取数据，查看缺失值情况
    \item 用中位数填充age列的缺失值
    \item 用平均值填充salary列的缺失值
    \item 用众数填充department列的缺失值
    \item 用众数填充gender列的缺失值
    \item 删除score列仍有缺失值的行
    \item 按department分组，统计平均工资和平均年龄
    \item 保存处理后的数据
\end{enumerate}

\textbf{参考答案：}
\begin{lstlisting}
import pandas as pd

# 读取数据
df = pd.read_csv('03_缺失值处理_练习.csv')

# 查看缺失值
print("缺失值统计：")
print(df.isnull().sum())

# 用中位数填充年龄
median_age = df['age'].median()
df['age'] = df['age'].fillna(median_age)

# 用平均值填充工资
mean_salary = df['salary'].mean()
df['salary'] = df['salary'].fillna(mean_salary)

# 用众数填充部门和性别
mode_dept = df['department'].mode()[0]
df['department'] = df['department'].fillna(mode_dept)

mode_gender = df['gender'].mode()[0]
df['gender'] = df['gender'].fillna(mode_gender)

# 删除score缺失的行
df = df.dropna(subset=['score'])

# 按部门统计
dept_stats = df.groupby('department').agg({
    'salary': 'mean',
    'age': 'mean'
}).round(2)

print("\n部门统计：")
print(dept_stats)

# 保存结果
df.to_csv('处理后的员工数据.csv', index=False)
\end{lstlisting}
\end{practicebox}

%========================================
\section{代码速查表}
%========================================

\subsection{文件操作}

\begin{lstlisting}[caption={读取和保存CSV}]
# 读取CSV
df = pd.read_csv('文件名.csv')

# 保存CSV
df.to_csv('新文件名.csv', index=False)
\end{lstlisting}

\subsection{数据查看}

\begin{lstlisting}[caption={数据查看}]
df.shape          # 形状（行数，列数）
df.columns        # 列名
df.head()         # 前5行
df.tail()         # 后5行
df.info()         # 详细信息
df.describe()     # 统计摘要
df.dtypes         # 数据类型
\end{lstlisting}

\subsection{数据选择}

\begin{lstlisting}[caption={数据选择}]
# 选择列
df['列名']
df[['列1', '列2']]

# 选择行
df.iloc[0]        # 第1行
df.iloc[0:5]      # 前5行

# 条件筛选
df[df['列名'] > 值]
\end{lstlisting}

\subsection{数据清洗}

\begin{lstlisting}[caption={数据清洗}]
# 缺失值
df.isnull().sum()                 # 统计缺失值
df['列'].fillna(值)              # 填充缺失值
df.dropna()                       # 删除缺失值

# 重复值
df.duplicated().sum()             # 统计重复值
df.drop_duplicates()              # 删除重复值

# 异常值（IQR方法）
Q1 = df['列'].quantile(0.25)
Q3 = df['列'].quantile(0.75)
IQR = Q3 - Q1
outliers = df[(df['列'] < Q1-1.5*IQR) | (df['列'] > Q3+1.5*IQR)]
\end{lstlisting}

\subsection{数据转换}

\begin{lstlisting}[caption={数据转换}]
# 类型转换
df['列'].astype(float)
df['列'].astype(str)

# 字符串处理
df['列'].str.replace('旧', '新')
df['列'].str.strip()

# 应用函数
df['新列'] = df['列'].apply(函数)
\end{lstlisting}

\subsection{分组统计}

\begin{lstlisting}[caption={分组统计}]
# 简单分组
df.groupby('分组列')['统计列'].mean()

# 多统计量
df.groupby('分组列').agg({
    '列1': 'mean',
    '列2': ['mean', 'sum']
})
\end{lstlisting}

\end{document}
